\chapter{Параметры роботизированной складской среды}\label{app:А}

\begin{table}[h]
	\centering
	\caption{Параметры AMR, используемые для калибровки симуляции}
	\label{tab:amr_real_params}
	\begin{tabular}{|p{0.20\textwidth}|p{0.22\textwidth}|p{0.45\textwidth}|}
		\hline
		\textbf{Параметр} &
		\textbf{Значение в симуляции} &
		\textbf{Обоснование} \\
		\hline
		Тип робота &
		Малый/средний складской AMR &
		Соответствует классу OTTO 100 и MiR250
		\cite{OTTO100Spec2026,MiR250Spec2026} \\
		\hline
		Максимальная скорость &
		2.0 м/с &
		В спецификациях и описаниях класса указывается
		скорость порядка 2.0 м/с \\
		\hline
		Рабочая скорость в модели &
		0.5--1.5 м/с &
		Модельное ограничение ниже максимальной скорости, учитывающее движение в
		проходах, повороты, зоны безопасности и взаимодействие с другими роботами \\
		\hline
		Полезная нагрузка &
		100--150 кг &
		Консервативная калибровка под малые и средние складские AMR; верхняя граница согласована с классом OTTO~100, а диапазон ниже паспортного максимума MiR250 \\
		\hline
		Время автономной работы &
		6--17.5 часов &
		Диапазон соответствует спецификациям OTTO 100 и MiR250 \\
		\hline
		Типовые датчики &
		Лидар, камеры, инерциальные датчики, датчики безопасности &
		Типичный набор сенсоров AMR в складской робототехнике
		\cite{Keith2024WarehouseAMRReview} \\
		\hline
		Сеть &
		Wi-Fi 2.4/5 ГГц; в стресс-сценариях моделируется ухудшенная связь &
		Для складского AMR-сценария предполагается локальная беспроводная сеть; в
		стресс-сценариях используются параметры задержки и пропускной способности из
		работ по offloading мобильных роботов \cite{Baruffa2025AMRResourceAssignment} \\
		\hline
	\end{tabular}
\end{table}

\begin{table}[h]
	\centering
	\caption{Ключевые параметры экспериментальной среды и операционных режимов}
	\label{tab:appendix_env_params}
	\begin{tabular}{|p{0.24\textwidth}|p{0.19\textwidth}|p{0.40\textwidth}|}
		\hline
		\textbf{Параметр} &
		\textbf{Значение} &
		\textbf{Назначение} \\
		\hline
		Длительность управляющего слота &
		100 мс &
		Единица дискретизации принятия решений в экспериментальной среде \\
		\hline
		Queue horizon &
		350 мс &
		Горизонт учёта очереди и отложенного ресурсного потребления \\
		\hline
		Queue reservation factor &
		0.2 &
		Доля ресурса, резервируемая при оценке допустимости назначения \\
		\hline
		Fixed normal exposure &
		0.55 &
		Фиксированный базовый режим для методов Fixed-Heuristic и Fixed-Auction \\
		\hline
		Safe exposure &
		0.30 &
		Защитный режим с минимальным раскрытием ёмкости узла \\
		\hline
		Aggressive exposure &
		0.90 &
		Режим максимального использования доступных ресурсов \\
		\hline
		Reserve exposure &
		0.45 &
		Режим резервирования под будущие критические задачи \\
		\hline
		Critical exposure &
		0.70 &
		Режим приоритизации классов \textit{obstacle detection}, \textit{localization/slam} и \textit{path replanning} \\
		\hline
		Peak scenario &
		$0.20$, spike $\times 4$, 36 шагов &
		Параметры сценария кратковременного всплеска нагрузки \\
		\hline
	\end{tabular}
\end{table}
